%!TEX root = ../thesis.tex
\chapter{Introduction}
\label{ch:intro}
\todo{sämtliche Abkürzungen mit glocaries einbinden und verwalten}
\section{Motivation}


The increasing availability of high-resolution satellite and aerial image data has made the global detection of objects technically possible. At the same time, the cost of the necessary computing power is falling, making it increasingly easier to use \acrfull{AI} methods for the fully automated analysis of such data \cite{Lemley2017}. This opens up new possibilities for the efficient evaluation of large amounts of image data, which are of considerable relevance in both civil and military contexts \cite{Lemley2017}.
 
One important area of application is the analysis and documentation of destroyed infrastructure or building structures in crisis and disaster areas, especially when direct on-site surveys are not possible due to security risks \cite{Knoth2017}. AI-supported image analysis can provide valuable support here, as aerial and satellite images enable efficient and safe information gathering.

Armed conflicts severely affect food security and the sustainable management of land and water resources. 
The destruction of infrastructure, population displacement, and the loss of agricultural areas lead to complex socio-ecological crises that threaten local livelihoods and have long-term global implications \cite{Mhanna2023}. 

Scientific analysis of such regions is complicated by the lack of \acrfull{GT} data, as field measurements are often impossible due to security concerns — as exemplified by the Orontes River Basin in Syria \cite{Mhanna2023}. 
To address this, multitemporal satellite analyses have been used to monitor \acrfull{LULC} changes, revealing significant declines in agricultural land and the emergence of refugee settlements \cite{Mhanna2023}. 

Recent studies have extended these approaches to map environmental impacts and so-called \emph{land scars} caused by armed conflicts using object-based image analysis (\acrshort{OBIA}) and deep convolutional neural networks (\acrshortpl{DCNN}) \cite{Feizizadeh2025}. 
Another sensitive but important research direction involves the remote detection of single and mass graves, where hyperspectral imaging has proven effective in identifying spectral anomalies in vegetation and soil that indicate decomposition processes \cite{Kalacska2006,Leblanc2014}.


Another application scenario is the analysis of the mobility behaviour of the population, for example by recording vehicles in large public car parks. Automated detection and classification methods can be used to identify and count vehicles on satellite or aerial images, allowing conclusions to be drawn about the utilisation of parking spaces. This information can be used to determine the need for additional parking spaces or to evaluate the efficiency of existing infrastructure. In addition, it opens up the possibility of quantifying traffic volumes at regional or national level, provided that continuous vehicle detection is implemented \cite{Eikvil2009}. A further question is the differentiation between moving and stationary vehicles, which can provide additional insights into mobility patterns and traffic flows.

There are also a wide range of possible applications in the context of security. The automated evaluation of large amounts of image data can help to identify changes in sensitive regions at an early stage and improve the assessment of the situation. A key advantage is that AI models reduce the workload for human analysts by automatically detecting and marking potential objects or structures before final validation takes place.

The use cases described in the previous sections illustrate the potential of remote sensing in combination with \acrfull{DL} for humanitarian, environmental and forensic issues. The data collected can be used to provide evidence and track war crimes in order to hold those responsible to account and identify missing persons. In addition, the insights gained open up opportunities for early detection and prevention of conflicts. By analysing mobility patterns and changes in land use across large areas, potential areas of tension can be identified at an early stage and appropriate measures for conflict reduction can be derived.

However, a key challenge in applying machine learning to high-resolution remote sensing images is the limited availability of suitable training data and the fact that vehicles are relatively small in scale. The precise detection and classification of these small objects requires specialised deep learning models that can reliably capture even fine structures.

\Acrlong{DL} is particularly suitable for the analysis of remote sensing data, as it can capture complex spatial-semantic patterns in high-resolution images that are difficult to model using classical, feature-based methods. In particular, the combination of multispectral channels, such as \acrshort{NIR} or \acrshort{NDVI}, and the integration of spatial context allows objects to be recognised more precisely and distinguished from background structures. Another challenge is the choice of bounding box strategy: \Acrfullpl{abb} can be inaccurate for obliquely oriented or overlapping objects, while \Acrfullpl{obb} enable more precise localisation and classification of small objects. This capability is crucial for reliable detection and classification of small objects in heterogeneous landscapes, such as those found in aerial or satellite images.

Preliminary work in this field, as carried out in a previous bachelor's thesis \cite{Balzer2022}, shows that deep learning methods such as \acrfull{YOLO} (v3) can generally detect small objects in simple scenarios. Through targeted adjustment of hyperparameters, the \acrfull{mAP} for very small objects could be improved. However, in the case of ship detection, it became apparent that the \acrshort{mAP} values achieved must be viewed critically and that the suitability of \acrshort{YOLO}v3 for the reliable detection of small objects is limited \cite{Balzer2022}.

It should be noted that the original \acrfull{YOLO} architecture was not specifically designed for the analysis of remote sensing data, but primarily for object recognition in everyday scenes, such as those found in conventional image datasets (e.g. \acrshort{COCO}\cite{lin2015} or ImageNet \cite{imagenet1,imagenet2}) \cite{redmon2016,yolov3_redmon,wang2024,yolov12}. Accordingly, the model is based on a different spatial-sensory context that differs significantly from remote sensing applications. This so-called \emph{domain shift} problem limits \acrshort{YOLO}'s performance in detecting small objects viewed from a bird's-eye perspective. These models are therefore optimised for the detection of large, clearly defined objects in a typical ground perspective. In the context of remote sensing, however, the conditions are fundamentally different: objects appear significantly smaller due to the recording perspective, sometimes overlap and are embedded in heterogeneous, complex landscapes.
 
The resulting methodological challenge is that classic detection models such as \acrshort{YOLO} have so far hardly taken the spatial context of objects into account. A central feature of this work is therefore the explicit integration of the spatial context into object detection. Real objects always occur in real, spatially structured environments whose background – for example, asphalt, sand or vegetation – significantly influences detection performance. The integration of such contextual information (e.g. through \acrshort{NDVI}) into the analysis is a key methodological focus of this work, as it contributes significantly to improving recognition accuracy in realistic remote sensing scenarios.

This master's thesis builds on these results and applies the question to the detection and classification of small vehicles in high-resolution multispectral aerial image data. A newer version of \acrshort{YOLO} (v9) is used, which features improvements in the network architecture. In addition, the extent to which the integration of true color and multispectral aerial images can increase the precision of the detection and classification of very small objects is investigated. In this way, the work expands existing research both in terms of content (focus on vehicles) and methodology (use of modern \acrshort{YOLO} architecture and multispectral data) and addresses the research gap in object recognition in complex scenarios. The extension of the analysis to include the spatial context – in particular by taking into account different background types such as asphalt or sand – enables a more realistic evaluation of model performance and thus addresses a central challenge of object recognition in remote sensing images.



\section{Structure of the thesis}


This thesis provides a comprehensive overview of the methods, experiments and results of the research. It begins with an introduction that outlines the motivation and objectives of the work, followed by the fundamentals, which explain the theoretical background and relevant metrics. The state of research summarises the current literature before the methodology describes the data sets, tools and procedures used in detail. The results of the experiments are then presented and interpreted in the discussion and placed in the context of existing research. The work concludes with a conclusion and an outlook on possible future research directions. Supplementary material such as images and tables can be found in the appendix.
